% Lösungen Einheit 2 von 3 – Kettenregel (zusammenbau v0.1)
\einheitenkopf{Einheit 2 von 3 · Kettenregel}

\erg{14}{a) der Exponent $7x$ \quad b) $f'(x) = 5 \cdot e^{5x}$ \quad c) $f'(x) = 4 \cdot 3 \cdot (3x + 1)^3 = 12 \cdot (3x + 1)^3$ \quad d) $f'(x) = 4x \cdot e^{2x^2}$ \quad e) $f^{(5)}(x) = 4^5 \cdot e^{4x}$; aus $e^{4(x + d)} = 4^5 \cdot e^{4x}$ folgt $e^{4d} = 4^5$, also $d = \frac{5}{4} \cdot \mathrm{ln}\,4 \approx 1{,}73$ – Verschiebung um etwa $1{,}73$ nach links. \quad f) $f(2) = 1$ und $g'(1) = 0$ (Tiefpunkt), also $h'(2) = g'(f(2)) \cdot f'(2) = 0$; $h(2) = g(1) = -2$; Tangente: $y = -2$}
\erg{15}{a) $f^{(5)}(x) = 2^5 \cdot \mathrm{cos}(2x) = 32 \cdot \mathrm{cos}(2x)$}
\erg{16}{a) Die innere Ableitung $-6$ fehlt; richtig: $f'(x) = -6 \cdot e^{-6x}$. \quad b) Die innere Funktion $4x$ hat die Ableitung $4$; nach der Kettenregel wird die äußere Ableitung $e^{4x}$ mit dieser inneren Ableitung multipliziert.}
